\section{Introduction}

\label{sec:intro}

The deployment of large language models (LLMs) in production environments demands inference frameworks that are simultaneously fast, memory-efficient, predictable, and safe. The dominant inference frameworks---PyTorch~\cite{paszke2019pytorch}, ONNX Runtime~\cite{onnxruntime2019}, TensorRT~\cite{tensorrt2017}---each make trade-offs that leave significant performance on the table.

\subsection{Limitations of Existing Frameworks}

\begin{enumerate}[leftmargin=1.4em]
    \item \textbf{Python runtime overhead}: PyTorch inference runs through the Python interpreter, incurring function call overhead, GIL contention for multi-threaded serving, and unpredictable garbage collection pauses that cause latency spikes.
    \item \textbf{Memory management}: PyTorch's reference-counted memory management leads to fragmentation and unpredictable allocation patterns during inference, particularly for variable-length sequence processing.
    \item \textbf{Kernel launch overhead}: Transformer models consist of hundreds of small operations (matrix multiplications, layer norms, activations), each requiring a separate kernel launch. The cumulative overhead is significant for small batch sizes typical of interactive inference.
    \item \textbf{Type safety}: Dimension mismatches, dtype errors, and device placement bugs are caught only at runtime, often deep in model execution, making debugging difficult and increasing the risk of production failures.
\end{enumerate}

\subsection{Why Rust?}

Rust~\cite{klabnik2019rust} provides three properties critical for inference frameworks:

\begin{itemize}[leftmargin=1.1em]
    \item \textbf{Zero-cost abstractions}: High-level APIs compile to the same machine code as hand-written C, with no runtime overhead.
    \item \textbf{Ownership system}: Compile-time memory management eliminates garbage collection pauses, ensuring predictable latency.
    \item \textbf{Type system}: Algebraic data types and trait-based generics enable compile-time verification of tensor operations.
\end{itemize}

\subsection{Contributions}

\begin{enumerate}[leftmargin=1.4em]
    \item A zero-copy tensor engine with arena allocation (\S\ref{sec:tensor}).
    \item A fused kernel compiler for transformer operations (\S\ref{sec:fusion}).
    \item A type-safe model definition API (\S\ref{sec:api}).
    \item Comprehensive benchmarks against PyTorch, ONNX Runtime, and llama.cpp (\S\ref{sec:evaluation}).
    \item Production deployment analysis (\S\ref{sec:production}).
\end{enumerate}

